\documentclass[11pt]{article}

% Standard, arXiv-safe packages only (keeps compile risk low).
\usepackage[utf8]{inputenc}
\usepackage[T1]{fontenc}
\usepackage{lmodern}
\usepackage[margin=1in]{geometry}
\usepackage{microtype}
\usepackage{booktabs}
\usepackage{array}
\usepackage{enumitem}
\usepackage{hyperref}
\hypersetup{colorlinks=true, linkcolor=black, citecolor=black, urlcolor=blue}
\usepackage{url}

\title{\textbf{falsify-ml}: A Falsification Protocol for Deciding\\
Whether an ML Project Is Worth Pursuing}
\author{Saurabh Dixit}
\date{2026}

\begin{document}
\maketitle

\begin{abstract}
\noindent\textit{Preprint / position paper. This is a methods-and-position paper, not an empirical
results paper. It presents a protocol and an open-source tool, situates them against existing
reproducibility and reporting work, and states honestly what it does \emph{not} yet establish (an
empirical evaluation of the protocol's hit-rate on known failures --- see Section~\ref{sec:limits}).
It has not been peer-reviewed.}

\medskip
\noindent Most machine-learning tooling helps practitioners \emph{build} models; comparatively little
helps them decide whether a model --- or the idea behind it --- is worth building or worth trusting
in the first place. Yet a large fraction of ML failures are decided \emph{before} a single training
run: a wrongly framed problem, a leaked train/test split, data that cannot physically support the
target precision, no baseline, or economics that never close. We present \textbf{falsify-ml}, a
ten-stage, model-agnostic protocol that treats every load-bearing ML claim as a \emph{hypothesis to
be falsified} rather than a result to be confirmed, and an accompanying open-source tool (pure Python
standard library, zero dependencies) that runs the mechanical parts of the protocol and produces a
single project-level verdict: \textbf{PURSUE / DO NOT PURSUE / INSUFFICIENT EVIDENCE}. The protocol
synthesizes ideas from the philosophy of science (Popper's falsification), pre-registration, and the
recent ML reproducibility literature (data-leakage detection; reporting standards) into a
\emph{pre-build and trust-gate} workflow with runnable diagnostics for the three stages where
automation has the most leverage: leakage/split-hygiene, data-feasibility (statistical power and
label-noise ceilings), and value/economics. We position falsify-ml relative to prior work, describe
the protocol and tool, work through illustrative cases, and are explicit about what we do \emph{not}
claim. Code: \url{https://github.com/saurabhdxt-labs/falsify-ml} (Apache-2.0).
\end{abstract}

\section{Introduction}

A model that scores beautifully on its own benchmark can be detecting nothing real. When the
benchmark and the training data are drawn from the same generator, an internal metric measures
\emph{consistency with the generator}, not agreement with reality --- a closed loop that produces a
confident number with no external meaning. This failure mode is neither rare nor new: Kapoor and
Narayanan~\cite{kapoor2023leakage} document data leakage across 17 scientific fields affecting 294
papers, and show that when leakage is corrected, the apparent superiority of ML methods frequently
disappears.

The expensive property of these failures is \emph{when} they are discovered. A leaked split, an
unachievable target given label noise, or a model whose work is actually done by a fallback path can
survive months of development and downstream building before anyone measures the claim against
independent evidence. The cost is not a bug fix; it is the wasted quarter.

The existing response has largely been \emph{reporting and review standards} applied to
\emph{finished} work --- most prominently REFORMS~\cite{kapoor2024reforms}, a consensus 32-item
checklist for what an ML-based-science paper should disclose. These are valuable, but they are
oriented toward catching problems at review time and toward reporting, not toward \emph{deciding,
before you build, whether the project can work at all}, and not toward \emph{adversarially attacking
your own claim} with runnable diagnostics.

falsify-ml addresses that gap. Its premises are:

\begin{enumerate}[leftmargin=*]
\item \textbf{A claim is a hypothesis until it is closed.} Every load-bearing assertion (``the model
works,'' ``this beats the baseline,'' ``the detector detects the thing'') is treated as a hypothesis
requiring a precise statement, a pre-registered TRUE shape and FALSE shape, evidence from
\emph{outside} the loop that produced the model, and a written verdict --- before downstream work
depends on it.
\item \textbf{Most failures are upstream of evaluation.} Falsification of the model's headline claim
is one stage of ten; the others guard problem framing, data hygiene, data feasibility, baselines,
economics, training health, eval design, deployment drift, and stopping.
\item \textbf{The useful output is a decision, not an analysis.} The protocol rolls per-stage statuses
into one verdict --- PURSUE / DO NOT PURSUE / INSUFFICIENT EVIDENCE --- usable on a bare idea (before
data exists) or on a real model with data.
\end{enumerate}

We make no claim that the individual techniques are novel (they are not; see
Section~\ref{sec:related}). The contribution is the \emph{synthesis} into a pre-build-and-trust
lifecycle with a runnable, dependency-free tool, and an explicit calibration that keeps the protocol
from degenerating into ``never ship anything.''

\section{The core idea: hypothesis closure}

A hypothesis is \emph{closed} when it has all six of: (i) a precise one-sentence statement (one
variable, one measurable outcome); (ii) a predicted TRUE shape, written before measurement; (iii) a
predicted FALSE shape (the counterfactual --- what you would see if the claim is wrong; if you cannot
write it, the claim is unfalsifiable and should be discarded); (iv) independent external evidence
measured against both shapes; (v) a verdict of WORKING (with named coverage) or NOT WORKING (with a
causal reason); (vi) the verdict written down and dated before downstream work depends on it. The
default verdict under ambiguity is NOT WORKING: the burden of proof is on the claim. ``Probably,''
``it seems to,'' ``partially,'' and ``we'll measure that later'' are not verdicts.

This is Popperian falsification~\cite{popper1959} applied operationally to ML claims, combined with
\emph{pre-registration} --- committing the success criteria before seeing results, a practice
imported from clinical trials and psychology to combat post-hoc rationalization, and increasingly
discussed for ML~\cite{prereg2020}. The novelty here is not the principle but the operational
packaging: a date-locked threshold file, a planted-signal test of the evaluation harness itself, and
a closed-loop diagnostic (below).

\section{The ten-stage lifecycle}

Falsification of the headline claim is \textbf{Stage 5}. The full set:

\begin{center}
\renewcommand{\arraystretch}{1.2}
\begin{tabular}{@{}c l p{6.2cm} l@{}}
\toprule
\# & Stage & Failure it catches & Tool \\
\midrule
0 & Frame the problem & building the wrong thing; using ML where a rule suffices & --- \\
1 & Data contract + leakage & train/test contamination & \texttt{leakage} \\
1.5 & Data reality & signal/precision the data cannot physically support & \texttt{reality} \\
2 & Baseline before training & inventing the bar after seeing the result & --- \\
2.5 & Value / economics & scientifically sound but economically pointless & \texttt{value} \\
3 & Training health & a model converged to noise & --- \\
4 & Eval design & a mis-chosen metric; underpowered evaluation & --- \\
5 & Falsify the claim & a claim that fits in-sample while measuring nothing real & \texttt{planted}, \texttt{closed-loop}, \texttt{prereg} \\
6 & Deploy + monitor & silent drift after deployment & --- \\
7 & Stopping rule & iterating forever; never revisiting a stale claim & --- \\
\bottomrule
\end{tabular}
\end{center}

Each stage emits \textbf{PASS / FAIL / UNKNOWN}. The roll-up rule is mechanical: any FAIL on a
critical stage (0, 1, 1.5, 2.5, 5) $\rightarrow$ \textbf{DO NOT PURSUE}; else any UNKNOWN on a
critical stage $\rightarrow$ \textbf{INSUFFICIENT EVIDENCE}; all critical PASS $\rightarrow$
\textbf{PURSUE}. The idea-only path (no data yet) leaves data-dependent stages UNKNOWN and typically
yields INSUFFICIENT EVIDENCE with a concrete ``collect a pilot dataset next'' --- but can yield DO NOT
PURSUE outright when a judgment or feasibility stage fails (e.g.\ a precision target above the
label-noise ceiling).

\section{The diagnostics}

Three stages have the most automation leverage; the tool implements them as task-agnostic,
pure-stdlib CLIs (no numpy/scipy/pandas), supporting classification (AUROC), regression (Pearson),
and ranking (Spearman) via a \texttt{--metric} flag.

\paragraph{Leakage / split hygiene (\texttt{leakage}, Stage 1).} Detects duplicate/near-duplicate rows
spanning the train/test split, temporal overlap (test not strictly after train), and group/entity
leakage (an entity in both splits) as hard failures; a single feature that perfectly partitions the
label is flagged as a \emph{suspect for human review}, not an auto-failure (legitimate
downstream-consequence features exist).

\paragraph{Data reality / feasibility (\texttt{reality}, Stage 1.5).} In \emph{numbers mode} (no data
yet) it computes the minimum detectable correlation at a given sample size via the Fisher-$z$
transform, and the achievable-precision ceiling implied by a stated label-noise rate
($\text{ceiling} = 1 - \text{error\_rate}$), returning a verdict on whether the target is reachable
\emph{in principle}. In \emph{CSV mode} it reports sample size, class balance, per-column missingness,
and label-noise from a re-labeled gold subset. The inverse-normal quantiles use Acklam's
approximation; we verified the power computation is bit-accurate against \texttt{scipy.stats} at
authoring time (e.g.\ $n{=}20 \rightarrow$ minimum detectable $r \approx 0.591$; $n{=}100 \rightarrow
0.277$).

\paragraph{Value / economics (\texttt{value}, Stage 2.5).} Computes gross gain, annual net, and
payback period from value-per-unit-of-improvement, expected lift, deploy cost, and maintenance cost,
returning PASS / FAIL / UNKNOWN. This catches the ``1\% lift, \$2M to deploy'' project that is
scientifically fine and economically dead.

\paragraph{Falsification harness (Stage 5).} Two diagnostics support the trust gate. The
\emph{planted-signal meta-falsification} (\texttt{planted}) injects an obvious signal and confirms the
chosen metric is near-perfect, then injects pure noise and confirms the metric collapses to its
no-skill floor --- a test of \emph{the evaluation machinery itself} before any verdict from it is
trusted. The \emph{closed-loop / external-oracle diagnostic} (\texttt{closed-loop}) correlates the
model's in-sample score against an \emph{independent} oracle on the same cases (joined by ID, never by
row order), with a percentile-bootstrap confidence interval, and flags the signature of a closed loop:
high in-sample fit paired with near-zero external correlation. A pre-registration utility
(\texttt{prereg}) date-locks the TRUE/FALSE shapes and thresholds with a content hash so post-hoc
threshold-softening is visible.

The tool ships with a test suite (119 tests at time of writing) in which statistical kernels are
checked against hand-derived or scipy-cross-checked values, and each test documents the code mutation
it is designed to catch.

\section{Illustrative use}
\label{sec:illustrative}

We give two short illustrative walkthroughs (not an evaluation --- see Section~\ref{sec:limits}).

\paragraph{Idea-only, no data --- employee-resignation prediction.} A practitioner wants $\geq 90\%$
precision but estimates labels are $\sim 15\%$ noisy. \texttt{reality --label-error-rate 0.15
--required-precision 0.90} returns DO NOT PURSUE: the precision ceiling is 0.85, below the
requirement, unreachable regardless of model quality --- even though statistical power (a 150-case
pilot can detect $r \geq 0.3$) and the economics both pass. The decision is reached before any data
is collected.

\paragraph{Data-stage --- a churn model with a planted leak.} On a synthetic 480-row dataset where 30
customers appear in both train and test, \texttt{leakage --group-col customer\_id} flags the group
leak and returns FAIL (exit non-zero), while \texttt{closed-loop} against an independent human-label
oracle shows the score \emph{does} carry real signal ($\rho \approx 0.37$, CI clears zero). The
composed verdict is DO NOT PURSUE \emph{until the split is fixed}: the signal is plausibly real, but
the headline metric a stakeholder would see is inflated by the leak --- a fixable PURSUE-candidate,
not a dead project.

\section{Relation to prior work}
\label{sec:related}

Each ingredient has antecedents; the contribution is the synthesis and the pre-build framing.

\begin{itemize}[leftmargin=*]
\item \textbf{Reproducibility \& leakage.} Kapoor and Narayanan~\cite{kapoor2023leakage} establish the
scale and mechanism of leakage-driven irreproducibility. falsify-ml's Stage 1 operationalizes a subset
of these checks as a runnable pre-training gate.
\item \textbf{Reporting standards.} REFORMS~\cite{kapoor2024reforms} is the closest relative: a
consensus checklist for \emph{reporting} ML-based science. The differences are purpose and timing.
REFORMS asks ``what should a finished study disclose so a reviewer can judge it?''; falsify-ml asks
``should this project be pursued at all, and does its central claim survive an adversarial attack?''
--- applied \emph{before} and \emph{during} the work, with runnable diagnostics and a go/no-go verdict
rather than a disclosure checklist. The two are complementary.
\item \textbf{Falsification \& pre-registration.} Popper~\cite{popper1959} supplies the epistemic
stance; pre-registration for ML~\cite{prereg2020} supplies the commit-criteria-first mechanic.
falsify-ml packages both operationally.
\item \textbf{Power and statistics.} The data-feasibility stage applies standard power analysis and a
simple label-noise precision ceiling; these are textbook, not novel --- the contribution is making
them a \emph{day-zero gate} rather than a post-hoc rationalization.
\end{itemize}

To our knowledge, the specific packaging --- a single pre-build-to-deployment lifecycle that (a)
treats the project-level ``is this worth pursuing?'' question as the top-level output, (b) ships
runnable, dependency-free diagnostics for the highest-leverage stages, and (c) works on a bare idea
before data exists --- is not directly covered by the above. We make this as a \emph{positioning}
claim, not a priority claim; we have not done an exhaustive literature survey and welcome pointers to
closer work.

\section{Limitations and what this paper does \emph{not} claim}
\label{sec:limits}

In keeping with the protocol's own discipline, we state the boundaries explicitly.

\begin{enumerate}[leftmargin=*]
\item \textbf{No empirical evaluation of the protocol's efficacy.} We have not measured how often
falsify-ml would have flagged real, independently-documented ML failures (e.g.\ the
retracted/corrected papers catalogued in~\cite{kapoor2023leakage}). The illustrative cases in
Section~\ref{sec:illustrative} are anecdotes and one is synthetic. The honest verdict on ``this
protocol catches real failures at rate $X$'' is \textbf{INSUFFICIENT EVIDENCE}; a retrospective study
on a corpus of known-leaky papers is the natural next step.
\item \textbf{No adoption or usability data.} Whether practitioners find the protocol usable and act
on its verdicts is untested.
\item \textbf{The estimates are only as good as their inputs.} The feasibility stage computes correct
statistics on \emph{user-supplied} effect sizes and label-noise rates; a confident-looking DO NOT
PURSUE built on a guessed input is itself a hazard.
\item \textbf{Scope is ML claims with data, labels, and a metric} --- classification, regression,
ranking, detection, forecasting. It is \emph{not} an evaluator of LLM/agent/generative quality, which
needs different machinery.
\item \textbf{The diagnostics cover three of ten stages.} The remaining stages are structured
checklists, not automated checks; their reliability depends on the practitioner answering them
honestly.
\end{enumerate}

A paper for a falsification tool that overstated its own evidence would be self-refuting. We have
tried to hold this paper to the standard the tool enforces.

\section{Conclusion}

falsify-ml reframes ML validation from ``find evidence it works'' to ``try hard to prove it doesn't,
and let it survive only if it genuinely does,'' and extends that stance across the whole project
lifecycle so that the cheapest place to kill a doomed project --- before data, before compute --- is
where the decision is made. The tool is open source, dependency-free, and task-agnostic. The
protocol's value is, by its own rules, an open hypothesis until evaluated against a corpus of known
failures; that evaluation is the work we invite.

\begin{thebibliography}{9}
\bibitem{kapoor2023leakage} S.~Kapoor and A.~Narayanan. Leakage and the reproducibility crisis in
machine-learning-based science. \emph{Patterns}, 4(9):100804, 2023.
\url{https://doi.org/10.1016/j.patter.2023.100804}
\bibitem{kapoor2024reforms} S.~Kapoor et al. REFORMS: Consensus-based Recommendations for
Machine-learning-based Science. \emph{Science Advances}, 2024. arXiv:2308.07832.
\url{https://doi.org/10.1126/sciadv.adk3452}
\bibitem{popper1959} K.~Popper. \emph{The Logic of Scientific Discovery}. 1959.
\bibitem{prereg2020} L.~Bertinetto, J.~F.~Henriques, S.~Albanie, et al.\ (eds.).
\emph{Proceedings of the NeurIPS 2020/2021 Workshops on Pre-registration in Machine Learning}.
PMLR Volumes 148 (2020) and 181 (2021). \url{https://proceedings.mlr.press/v148/}. See also
J.~M.~Hofman, A.~Chatzimparmpas, A.~Sharma, D.~J.~Watts, and J.~Hullman, Pre-registration for
Predictive Modeling, arXiv:2311.18807, 2023.
\end{thebibliography}

\vspace{1em}
\noindent\textit{\small This paper is methodology and positioning only. It contains no proprietary
model, data, or result. The tool's diagnostics run on the user's own project.}

\end{document}
